\section{Introdução}
Este documento apresenta as especificações dos requisitos 
funcionais do sistema por meio de histórias de usuário. 
Cada história de usuário é escrita sob a perspectiva do usuário final,
em linguagem informal, e segue estritamente o formato: 
\textbf{``Como [papel] ... Eu quero [ação] ... Para [benefício/valor]''}.

\section{Lista de Histórias de Usuário}

\subsection{Gerenciamento de Quadro e Raias}
\begin{itemize}
    \item \textbf{HU01 - Criação de Quadro:} Como usuário, eu quero criar um novo 
    quadro Kanban, para organizar um novo projeto da minha equipe.
    \item \textbf{HU02 - Visualização de Quadro:} Como usuário, eu quero ler/visualizar
    o quadro Kanban, para acompanhar o status atual das atividades.
    \item \textbf{HU03 - Atualização de Quadro:} Como usuário, eu quero atualizar 
    o nome e configurações do quadro, para mantê-lo alinhado com o contexto do 
    projeto.
    \item \textbf{HU04 - Exclusão de Quadro:} Como usuário, eu quero excluir um 
    quadro Kanban, para remover projetos finalizados ou cancelados.
    
    \item \textbf{HU05 - Criação de Raia (Swimlane):} Como usuário, eu quero criar 
    uma raia (swimlane) no quadro, para separar atividades por categorias, 
    prioridades ou times.
    \item \textbf{HU06 - Visualização de Raia:} Como usuário, eu quero visualizar 
    as raias no quadro, para ver as atividades agrupadas de forma clara.
    \item \textbf{HU07 - Atualização de Raia:} Como usuário, eu quero atualizar as 
    propriedades de uma raia, para corrigir ou alterar sua classificação.
    \item \textbf{HU08 - Exclusão de Raia:} Como usuário, eu quero excluir uma raia, 
    para remover divisões que não são mais necessárias.
\end{itemize}

\subsection{Gerenciamento de Cartões e Fluxo}
\begin{itemize}
    \item \textbf{HU09 - Criação de Cartão:} Como usuário, eu quero criar um cartão
    de atividades (contendo ID, nome, responsável, data limite, prioridade e 
    descrição), para documentar uma nova tarefa.
    \item \textbf{HU10 - Leitura de Cartão:} Como usuário, eu quero ler os detalhes
    de um cartão, para entender o que deve ser feito na tarefa.
    \item \textbf{HU11 - Atualização de Cartão:} Como usuário, eu quero atualizar
    as informações de um cartão, para refletir mudanças ou corrigir dados.
    \item \textbf{HU12 - Exclusão de Cartão:} Como usuário, eu quero excluir um 
    cartão, para remover tarefas duplicadas ou descartadas.
    \item \textbf{HU13 - Movimentação entre Colunas:} Como usuário, eu quero 
    mover um cartão de atividades entre as colunas do quadro (ex: A FAZER para 
    FAZENDO), para indicar o progresso do trabalho.
    \item \textbf{HU14 - Movimentação entre Raias:} Como usuário, eu quero mover um
    cartão de atividades entre as raias do quadro, para reclassificar a atividade.
\end{itemize}

\subsection{Limites e Métricas}
\begin{itemize}
    \item \textbf{HU15 - Configuração de Limite WIP:} Como gestor/administrador do 
    quadro, eu quero estabelecer o número máximo de cartões por coluna (work-in-progress 
    limit), para evitar sobrecarga da equipe.
    \item \textbf{HU16 - Cálculo de Métricas (Cycle Time):} Como gestor, eu quero que
    o sistema calcule a métrica Cycle Time, para saber quanto tempo em média a equipe 
    leva para concluir uma tarefa após iniciada.
    \item \textbf{HU17 - Cálculo de Métricas (Lead Time):} Como gestor, eu quero que
    o sistema calcule a métrica Lead Time, para saber o tempo total desde a criação da 
    demanda até a entrega.
    \item \textbf{HU18 - Cálculo de Métricas (Throughput):} Como gestor, eu quero 
    que o sistema calcule o Throughput (taxa de entrega), para saber quantos itens 
    estão sendo entregues em um determinado período.
    \item \textbf{HU19 - Cálculo de Métricas (WIP Atual):} Como gestor, eu quero que
    o sistema calcule/mostre o WIP atual (Work In Progress), para saber quantos itens 
    estão sendo trabalhados ativamente no momento.
\end{itemize}

\subsection{Autenticação e Autorização}
\begin{itemize}
    \item \textbf{HU20 - Login:} Como usuário, eu quero fazer login no sistema com validação em banco de dados relacional, para 
    acessar os quadros Kanban.
    \item \textbf{HU21 - Logout:} Como usuário, eu quero fazer logout do sistema, para 
    encerrar a sessão.
    \item \textbf{HU22 - Cadastro:} Como usuário, eu quero me cadastrar no sistema, para 
    criar uma conta.
    \item \textbf{HU23 - Compartilhar projetos:} Como usuário, eu quero adicionar pessoas a um projeto, para 
    trabalhar em conjunto.
\end{itemize}
